\documentclass[12pt]{article}
\usepackage{amsmath,amssymb,amsfonts}
\usepackage[margin=1in]{geometry}

\begin{document}

\section*{Chapter 3, Problem 16}

\textbf{Problem.} Show that the eigenvalues of the matrix
\[
M = \begin{pmatrix} 5 & 4 \\ 1 & 2 \end{pmatrix}
\]
are $\lambda = 1, 6$, and find the corresponding eigenvectors.

Construct a diagonalizing matrix $P$, prove that its inverse exists, compute its inverse, and verify that $P^{-1}MP$ is diagonal with the eigenvalues on the diagonal.

Note that $\det M = 6$. Is it true that any diagonalizable matrix with an eigenvalue equal to zero is singular?

\bigskip
\textbf{Solution.}

\textbf{Eigenvalues:} The characteristic equation is
\[
\det(M - \lambda I) = \det\begin{pmatrix} 5-\lambda & 4 \\ 1 & 2-\lambda \end{pmatrix} = (5-\lambda)(2-\lambda) - 4 = 0.
\]
Expanding:
\[
\lambda^2 - 7\lambda + 10 - 4 = \lambda^2 - 7\lambda + 6 = (\lambda - 1)(\lambda - 6) = 0.
\]
So $\lambda_1 = 1$ and $\lambda_2 = 6$.

\medskip
\textbf{Eigenvectors:}

For $\lambda_1 = 1$: $(M - I)x = 0$:
\[
\begin{pmatrix} 4 & 4 \\ 1 & 1 \end{pmatrix}\begin{pmatrix} x_1 \\ x_2 \end{pmatrix} = 0 \quad \Rightarrow \quad x_1 + x_2 = 0 \quad \Rightarrow \quad x^{(1)} = \begin{pmatrix} 1 \\ -1 \end{pmatrix}.
\]

For $\lambda_2 = 6$: $(M - 6I)x = 0$:
\[
\begin{pmatrix} -1 & 4 \\ 1 & -4 \end{pmatrix}\begin{pmatrix} x_1 \\ x_2 \end{pmatrix} = 0 \quad \Rightarrow \quad x_1 = 4x_2 \quad \Rightarrow \quad x^{(2)} = \begin{pmatrix} 4 \\ 1 \end{pmatrix}.
\]

\medskip
\textbf{Diagonalization:} Form $P$ with eigenvectors as columns:
\[
P = \begin{pmatrix} 1 & 4 \\ -1 & 1 \end{pmatrix}.
\]
$\det P = (1)(1) - (4)(-1) = 1 + 4 = 5 \neq 0$, so $P^{-1}$ exists.

\[
P^{-1} = \frac{1}{5}\begin{pmatrix} 1 & -4 \\ 1 & 1 \end{pmatrix}.
\]

Verify:
\[
P^{-1}MP = \frac{1}{5}\begin{pmatrix} 1 & -4 \\ 1 & 1 \end{pmatrix}\begin{pmatrix} 5 & 4 \\ 1 & 2 \end{pmatrix}\begin{pmatrix} 1 & 4 \\ -1 & 1 \end{pmatrix}.
\]
First compute $P^{-1}M$:
\[
P^{-1}M = \frac{1}{5}\begin{pmatrix} 1 & -4 \\ 1 & 1 \end{pmatrix}\begin{pmatrix} 5 & 4 \\ 1 & 2 \end{pmatrix} = \frac{1}{5}\begin{pmatrix} 5-4 & 4-8 \\ 5+1 & 4+2 \end{pmatrix} = \frac{1}{5}\begin{pmatrix} 1 & -4 \\ 6 & 6 \end{pmatrix}.
\]
Then:
\[
P^{-1}MP = \frac{1}{5}\begin{pmatrix} 1 & -4 \\ 6 & 6 \end{pmatrix}\begin{pmatrix} 1 & 4 \\ -1 & 1 \end{pmatrix} = \frac{1}{5}\begin{pmatrix} 1+4 & 4-4 \\ 6-6 & 24+6 \end{pmatrix} = \frac{1}{5}\begin{pmatrix} 5 & 0 \\ 0 & 30 \end{pmatrix} = \begin{pmatrix} 1 & 0 \\ 0 & 6 \end{pmatrix}.
\]

This confirms $P^{-1}MP = \text{diag}(1, 6)$. $\checkmark$

\medskip
\textbf{On the final question:} Yes. If a diagonalizable matrix has an eigenvalue equal to zero, then $P^{-1}AP = D$ where $D$ has a zero on the diagonal. Therefore $\det(D) = 0$, and since $\det(A) = \det(D)$ (similarity preserves determinant), $\det(A) = 0$, so $A$ is singular. In fact, this is true for \emph{any} matrix (diagonalizable or not): $\lambda = 0$ is an eigenvalue iff $\det(A - 0 \cdot I) = \det(A) = 0$. \hfill $\blacksquare$

\end{document}
